% Generated by tools/rao_tuning_report.py.
\begin{table*}[t]\centering\footnotesize
\caption{Held-out parameter validation: four paired IMU/initialization seeds and eight
vessel records, 32 cases per family. Entries average per-record RMS.
Violations count unchanged executable thresholds, not failed records.}
\label{tab:rao-parameter-validation}
\begin{tabular}{@{}llrrrrrr@{}}\toprule
Family & Setting & Violations & Yaw [deg] & Roll [deg] & Pitch [deg] & 3-D [m] & Worst/limit \\
\midrule
OU--II & baseline & 84 & 2.644 & 0.278 & 0.420 & 0.393 & 8.783 \\
OU--II & selected & 86 & 2.302 & 0.273 & 0.391 & 0.375 & 6.909 \\
OU--III & baseline & 90 & 2.469 & 0.282 & 0.356 & 0.311 & 8.883 \\
OU--III & selected & 88 & 2.128 & 0.285 & 0.331 & 0.291 & 6.739 \\
TFG & baseline & 44 & 2.967 & 0.352 & 0.501 & 0.478 & 6.778 \\
TFG & selected & 32 & 2.215 & 0.263 & 0.402 & 0.292 & 4.840 \\
\bottomrule\end{tabular}\end{table*}
Selected covariance and calibration settings improve aggregate errors but do not
make every quality gate pass. OU--II's held-out violation count increases despite
lower mean errors; OU--III roll RMS increases slightly. PII's held-out vertical
violations decrease from 23 to 7, with mean vertical RMS decreasing from
8.375\% to 5.961\% of incident $H_s$. The validation draws change sensor and
initialization errors on the same wave records; they do not cover unseen hulls.
